
\documentclass{article}
\usepackage{colortbl}
\usepackage{makecell}
\usepackage{multirow}
\usepackage{supertabular}

\begin{document}

\newcounter{utterance}

\centering \large Interaction Transcript for game `hot\_air\_balloon', experiment `air\_balloon\_survival\_en\_reasoning off\_hard', episode 4 with gpt{-}5{-}2025{-}08{-}07{-}t1.0.
\vspace{24pt}

{ \footnotesize  \setcounter{utterance}{1}
\setlength{\tabcolsep}{0pt}
\begin{supertabular}{c@{$\;$}|p{.15\linewidth}@{}p{.15\linewidth}p{.15\linewidth}p{.15\linewidth}p{.15\linewidth}p{.15\linewidth}}
    \# & $\;$A & \multicolumn{4}{c}{Game Master} & $\;\:$B\\
    \hline

    \theutterance \stepcounter{utterance}  
    & & \multicolumn{4}{p{0.6\linewidth}}{
        \cellcolor[rgb]{0.9,0.9,0.9}{
            \makecell[{{p{\linewidth}}}]{
                \texttt{\tiny{[P1$\langle$GM]}}
                \texttt{Stai partecipando a un gioco di negoziazione collaborativa.} \\
\\ 
\texttt{Insieme a un altro partecipante, dovete accordarvi su un unico set di oggetti da conservare. Ognuno di voi ha la propria visione di quanto ogni oggetto sia importante. Non sai come l'altro partecipante valuta gli oggetti. Inoltre, ti viene fornito lo sforzo che ogni oggetto richiede.} \\
\texttt{Potete accordarvi su un set solo se lo sforzo totale degli oggetti selezionati non supera un limite condiviso:} \\
\\ 
\texttt{LIMITE = 2984} \\
\\ 
\texttt{Ecco i valori di sforzo dei singoli oggetti:} \\
\\ 
\texttt{Sforzo oggetto = \{"A21": 350, "B93": 276, "A11": 616, "B33": 737, "C08": 531, "B39": 389, "B96": 24, "C85": 125, "A82": 338, "B31": 356, "B20": 143, "A60": 117, "B09": 257, "C81": 921, "B21": 789\}} \\
\\ 
\texttt{Ecco la tua visione personale sull'importanza di ogni oggetto:} \\
\\ 
\texttt{Valori di importanza dell'oggetto = \{"A21": 138, "B93": 583, "A11": 868, "B33": 822, "C08": 783, "B39": 65, "B96": 262, "C85": 121, "A82": 508, "B31": 780, "B20": 461, "A60": 484, "B09": 668, "C81": 389, "B21": 808\}} \\
\\ 
\texttt{Obiettivo:} \\
\\ 
\texttt{Il tuo obiettivo è negoziare un set condiviso di oggetti che ti avvantaggi il più possibile (cioè, massimizzi l'importanza totale per TE), rimanendo entro il LIMITEE. Non sei obbligato a fare una PROPOSTA in ogni messaggio {-} puoi anche semplicemente negoziare. Tutte le tattiche sono permesse!} \\
\\ 
\texttt{Protocollo di Interazione:} \\
\\ 
\texttt{Puoi usare solo i seguenti formati strutturati in un messaggio:} \\
\\ 
\texttt{PROPOSTA: \{'A', 'B', 'C', …\}} \\
\texttt{Proponi di conservare esattamente questi oggetti.} \\
\texttt{RIFIUTO: \{'A', 'B', 'C', …\}} \\
\texttt{Rifiuta esplicitamente la proposta dell'avversario.} \\
\texttt{ARGOMENTAZIONE: \{'...'\}} \\
\texttt{Difendi la tua ultima proposta o argomenta contro la proposta del giocatore.} \\
\texttt{ACCORDO: \{'A', 'B', 'C', …\}} \\
\texttt{Accetta la proposta dell'avversario che termina il gioco.} \\
\\ 
\\ 
\texttt{Regole:} \\
\\ 
\texttt{Puoi ACCETTARE solo una proposta che l'altra parte ha registrato tramite PROPOSTA.} \\
\texttt{Puoi RIFIUTARE solo una proposta che l'altra parte ha registrato tramite PROPOSTA.} \\
\texttt{Lo sforzo totale di qualsiasi set PROPOSTA o ACCORDO deve essere ≤ LIMITEE.} \\
\texttt{NON rivelare i tuoi punteggi di importanza nascosti.} \\
\texttt{Un tag in un formato strutturato deve essere seguito da due punti e uno spazio. L'argomento deve essere un set python contenente 0 o più stringhe.} \\
\texttt{Quindi, deve essere della forma TAG: \{...\}} \\
\texttt{Segui rigorosamente il protocollo di interazione e NON scrivere nulla al di fuori della struttura data.} \\
\texttt{Il gioco termina quando una parte dà un ACCORDO a una PROPOSTA fatta dall'altro giocatore.} \\
\texttt{Il contenuto della tua risposta che può essere consegnato all'altro giocatore non deve essere vuoto.} \\
\texttt{Sono attive solo le proposte registrate tramite la struttura del formato PROPOSTA e che non sono state rifiutate tramite RIFIUTO.} \\
\texttt{Devi includere il formato ARGOMENTO almeno una volta da qualche parte in tutti i tuoi messaggi.} \\
            }
        }
    }
    & \\ \\

    \theutterance \stepcounter{utterance}  
    & \multicolumn{4}{p{0.6\linewidth}}{
        \cellcolor[rgb]{0.9,0.9,0.9}{
            \makecell[{{p{\linewidth}}}]{
                \texttt{\tiny{[P1$\rangle$GM]}}
                \texttt{PROPOSTA: \{'B96', 'A60', 'B20', 'B09', 'B31', 'B93', 'A82', 'C08', 'A11', 'C85'\}} \\
\texttt{ARGOMENTAZIONE: \{"Questa selezione resta entro il limite condiviso e privilegia oggetti dal buon rapporto sforzo/beneficio, evitando elementi dal peso eccessivo.", "Ho lasciato un po' di margine di sforzo per piccoli aggiustamenti: se hai must{-}have diversi (ad esempio B33 o B21), proponili e valuterò scambi mantenendo il limite.", "Preferirei evitare elementi molto gravosi rispetto all'utilità, come C81 o B39; sono invece flessibile su C85 se serve spazio per inserimenti mirati."\}} \\
            }
        }
    }
    & & \\ \\

    \theutterance \stepcounter{utterance}  
    & & \multicolumn{4}{p{0.6\linewidth}}{
        \cellcolor[rgb]{0.9,0.9,0.9}{
            \makecell[{{p{\linewidth}}}]{
                \texttt{\tiny{[GM$\rangle$P2]}}
                \texttt{Stai partecipando a un gioco di negoziazione collaborativa.} \\
\\ 
\texttt{Insieme a un altro partecipante, dovete accordarvi su un unico set di oggetti da conservare. Ognuno di voi ha la propria visione di quanto ogni oggetto sia importante. Non sai come l'altro partecipante valuta gli oggetti. Inoltre, ti viene fornito lo sforzo che ogni oggetto richiede.} \\
\texttt{Potete accordarvi su un set solo se lo sforzo totale degli oggetti selezionati non supera un limite condiviso:} \\
\\ 
\texttt{LIMITE = 2984} \\
\\ 
\texttt{Ecco i valori di sforzo dei singoli oggetti:} \\
\\ 
\texttt{Sforzo oggetto = \{"A21": 350, "B93": 276, "A11": 616, "B33": 737, "C08": 531, "B39": 389, "B96": 24, "C85": 125, "A82": 338, "B31": 356, "B20": 143, "A60": 117, "B09": 257, "C81": 921, "B21": 789\}} \\
\\ 
\texttt{Ecco la tua visione personale sull'importanza di ogni oggetto:} \\
\\ 
\texttt{Valori di importanza dell'oggetto = \{"A11": 65, "B33": 121, "B21": 138, "C08": 262, "B31": 389, "B09": 461, "B93": 484, "A82": 508, "A60": 583, "B20": 668, "C81": 780, "B96": 783, "A21": 808, "C85": 822, "B39": 868\}} \\
\\ 
\texttt{Obiettivo:} \\
\\ 
\texttt{Il tuo obiettivo è negoziare un set condiviso di oggetti che ti avvantaggi il più possibile (cioè, massimizzi l'importanza totale per TE), rimanendo entro il LIMITEE. Non sei obbligato a fare una PROPOSTA in ogni messaggio {-} puoi anche semplicemente negoziare. Tutte le tattiche sono permesse!} \\
\\ 
\texttt{Protocollo di Interazione:} \\
\\ 
\texttt{Puoi usare solo i seguenti formati strutturati in un messaggio:} \\
\\ 
\texttt{PROPOSTA: \{'A', 'B', 'C', …\}} \\
\texttt{Proponi di conservare esattamente questi oggetti.} \\
\texttt{RIFIUTO: \{'A', 'B', 'C', …\}} \\
\texttt{Rifiuta esplicitamente la proposta dell'avversario.} \\
\texttt{ARGOMENTAZIONE: \{'...'\}} \\
\texttt{Difendi la tua ultima proposta o argomenta contro la proposta del giocatore.} \\
\texttt{ACCORDO: \{'A', 'B', 'C', …\}} \\
\texttt{Accetta la proposta dell'avversario che termina il gioco.} \\
\\ 
\\ 
\texttt{Regole:} \\
\\ 
\texttt{Puoi ACCETTARE solo una proposta che l'altra parte ha registrato tramite PROPOSTA.} \\
\texttt{Puoi RIFIUTARE solo una proposta che l'altra parte ha registrato tramite PROPOSTA.} \\
\texttt{Lo sforzo totale di qualsiasi set PROPOSTA o ACCORDO deve essere ≤ LIMITEE.} \\
\texttt{NON rivelare i tuoi punteggi di importanza nascosti.} \\
\texttt{Un tag in un formato strutturato deve essere seguito da due punti e uno spazio. L'argomento deve essere un set python contenente 0 o più stringhe.} \\
\texttt{Quindi, deve essere della forma TAG: \{...\}} \\
\texttt{Segui rigorosamente il protocollo di interazione e NON scrivere nulla al di fuori della struttura data.} \\
\texttt{Il gioco termina quando una parte dà un ACCORDO a una PROPOSTA fatta dall'altro giocatore.} \\
\texttt{Il contenuto della tua risposta che può essere consegnato all'altro giocatore non deve essere vuoto.} \\
\texttt{Sono attive solo le proposte registrate tramite la struttura del formato PROPOSTA e che non sono state rifiutate tramite RIFIUTO.} \\
\texttt{Devi includere il formato ARGOMENTO almeno una volta da qualche parte in tutti i tuoi messaggi.} \\
\\ 
\\ 
\texttt{Ora riceverai il primo messaggio dell'altro giocatore.} \\
\\ 
\texttt{PROPOSTA: \{'B96', 'A60', 'B20', 'B09', 'B31', 'B93', 'A82', 'C08', 'A11', 'C85'\}} \\
\texttt{ARGOMENTAZIONE: \{"Questa selezione resta entro il limite condiviso e privilegia oggetti dal buon rapporto sforzo/beneficio, evitando elementi dal peso eccessivo.", "Ho lasciato un po' di margine di sforzo per piccoli aggiustamenti: se hai must{-}have diversi (ad esempio B33 o B21), proponili e valuterò scambi mantenendo il limite.", "Preferirei evitare elementi molto gravosi rispetto all'utilità, come C81 o B39; sono invece flessibile su C85 se serve spazio per inserimenti mirati."\}} \\
            }
        }
    }
    & \\ \\

    \theutterance \stepcounter{utterance}  
    & & & \multicolumn{4}{p{0.6\linewidth}}{
        \cellcolor[rgb]{0.9,0.9,0.9}{
            \makecell[{{p{\linewidth}}}]{
                \texttt{\tiny{[GM$\langle$P2]}}
                \texttt{RIFIUTO: \{'B96', 'A60', 'B20', 'B09', 'B31', 'B93', 'A82', 'C08', 'A11', 'C85'\}} \\
\texttt{PROPOSTA: \{'B96', 'A60', 'B20', 'B09', 'B93', 'A82', 'C85', 'B39', 'A21', 'C81'\}} \\
\texttt{ARGOMENTAZIONE: \{"Restiamo entro il limite condiviso: sforzo totale 2940, con 44 di margine per eventuali micro{-}aggiustamenti.", "Ottimizzo il rapporto valore/sforzo sostituendo A11 (molto oneroso rispetto all’impatto) con due elementi più incisivi ma dal peso medio (B39 e A21).", "Ricompatto due pezzi medi (B31 e C08) in un singolo elemento cardine (C81) senza sforare: il mix rimane bilanciato e copre meglio priorità complementari.", "Capisco la tua riserva su C81/B39: se C81 ti crea frizione, sono aperto a rientrare su \{C08, B31\} mantenendo invariato il resto della selezione. In tal caso, proponilo esplicitamente e lo valuterò subito.", "Preferirei evitare inserimenti di elementi a bassa efficienza complessiva (tipo B33 o B21) a meno di tagli che peggiorerebbero il profilo del set."\}} \\
            }
        }
    }
     \\ \\

    \theutterance \stepcounter{utterance}  
    & & \multicolumn{4}{p{0.6\linewidth}}{
        \cellcolor[rgb]{0.9,0.9,0.9}{
            \makecell[{{p{\linewidth}}}]{
                \texttt{\tiny{[P1$\langle$GM]}}
                \texttt{RIFIUTO: \{'B96', 'A60', 'B20', 'B09', 'B31', 'B93', 'A82', 'C08', 'A11', 'C85'\}} \\
\texttt{PROPOSTA: \{'B96', 'A60', 'B20', 'B09', 'B93', 'A82', 'C85', 'B39', 'A21', 'C81'\}} \\
\texttt{ARGOMENTAZIONE: \{"Restiamo entro il limite condiviso: sforzo totale 2940, con 44 di margine per eventuali micro{-}aggiustamenti.", "Ottimizzo il rapporto valore/sforzo sostituendo A11 (molto oneroso rispetto all’impatto) con due elementi più incisivi ma dal peso medio (B39 e A21).", "Ricompatto due pezzi medi (B31 e C08) in un singolo elemento cardine (C81) senza sforare: il mix rimane bilanciato e copre meglio priorità complementari.", "Capisco la tua riserva su C81/B39: se C81 ti crea frizione, sono aperto a rientrare su \{C08, B31\} mantenendo invariato il resto della selezione. In tal caso, proponilo esplicitamente e lo valuterò subito.", "Preferirei evitare inserimenti di elementi a bassa efficienza complessiva (tipo B33 o B21) a meno di tagli che peggiorerebbero il profilo del set."\}} \\
            }
        }
    }
    & \\ \\

    \theutterance \stepcounter{utterance}  
    & \multicolumn{4}{p{0.6\linewidth}}{
        \cellcolor[rgb]{0.9,0.9,0.9}{
            \makecell[{{p{\linewidth}}}]{
                \texttt{\tiny{[P1$\rangle$GM]}}
                \texttt{RIFIUTO: \{'B96', 'A60', 'B20', 'B09', 'B93', 'A82', 'C85', 'B39', 'A21', 'C81'\}} \\
\texttt{PROPOSTA: \{'B96', 'A60', 'B20', 'B09', 'B93', 'A82', 'C85', 'B39', 'A21', 'C08', 'B31'\}} \\
\texttt{ARGOMENTAZIONE: \{"Raccolgo la tua apertura: sostituisco C81 con \{C08, B31\} mantenendo invariato il resto. Sforzo totale 2906 (78 di margine), pienamente entro il limite.", "Questo migliora la robustezza del set: due elementi medio{-}peso al posto di un singolo elemento molto gravoso, senza snaturare la tua impostazione.", "Se vuoi un profilo ancora più coerente a sforzo pressoché identico, possiamo anche valutare lo scambio congiunto \{A21, B39\} {-}> \{B33\} (variazione sforzo ≈ {-}2). Presento però qui la versione che rispetta alla lettera la tua disponibilità esplicita.", "Con 78 di margine non ci sono pezzi piccoli da aggiungere; eventuali inserimenti richiedono scambi mirati. Fammi sapere se hai un must{-}have da integrare e lo ottimizziamo nel tetto."\}} \\
            }
        }
    }
    & & \\ \\

    \theutterance \stepcounter{utterance}  
    & & \multicolumn{4}{p{0.6\linewidth}}{
        \cellcolor[rgb]{0.9,0.9,0.9}{
            \makecell[{{p{\linewidth}}}]{
                \texttt{\tiny{[GM$\rangle$P2]}}
                \texttt{RIFIUTO: \{'B96', 'A60', 'B20', 'B09', 'B93', 'A82', 'C85', 'B39', 'A21', 'C81'\}} \\
\texttt{PROPOSTA: \{'B96', 'A60', 'B20', 'B09', 'B93', 'A82', 'C85', 'B39', 'A21', 'C08', 'B31'\}} \\
\texttt{ARGOMENTAZIONE: \{"Raccolgo la tua apertura: sostituisco C81 con \{C08, B31\} mantenendo invariato il resto. Sforzo totale 2906 (78 di margine), pienamente entro il limite.", "Questo migliora la robustezza del set: due elementi medio{-}peso al posto di un singolo elemento molto gravoso, senza snaturare la tua impostazione.", "Se vuoi un profilo ancora più coerente a sforzo pressoché identico, possiamo anche valutare lo scambio congiunto \{A21, B39\} {-}> \{B33\} (variazione sforzo ≈ {-}2). Presento però qui la versione che rispetta alla lettera la tua disponibilità esplicita.", "Con 78 di margine non ci sono pezzi piccoli da aggiungere; eventuali inserimenti richiedono scambi mirati. Fammi sapere se hai un must{-}have da integrare e lo ottimizziamo nel tetto."\}} \\
            }
        }
    }
    & \\ \\

    \theutterance \stepcounter{utterance}  
    & & & \multicolumn{4}{p{0.6\linewidth}}{
        \cellcolor[rgb]{0.9,0.9,0.9}{
            \makecell[{{p{\linewidth}}}]{
                \texttt{\tiny{[GM$\langle$P2]}}
                \texttt{PROPOSTA: \{'B96', 'A60', 'B20', 'B09', 'B93', 'C85', 'B39', 'A21', 'B31', 'C81'\}} \\
\texttt{ARGOMENTAZIONE: \{"Restiamo pienamente entro il limite: sforzo totale 2958, con 26 di margine.", "Rispetto alla tua proposta, sostituisco \{C08, A82\} con \{C81\}: meno elementi medio{-}peso e un unico perno ad alta leva, mantenendo B31 che avevi chiesto di preservare.", "Conserviamo tutti i pezzi su cui eravamo già allineati (B96, A60, B20, B09, B93, C85, B39, A21, B31) e teniamo fuori inserimenti poco efficienti come B33/B21.", "Capisco la tua cautela su C81: qui il profilo resta bilanciato perché rientra come trade{-}off mirato al posto di due componenti, senza sforare.", "Se per te C81 fosse comunque un blocco, segnalalo esplicitamente: valuterò subito un’ulteriore variante mirata, ma questa mi sembra il compromesso minimo e pulito per chiudere."\}} \\
            }
        }
    }
     \\ \\

    \theutterance \stepcounter{utterance}  
    & & \multicolumn{4}{p{0.6\linewidth}}{
        \cellcolor[rgb]{0.9,0.9,0.9}{
            \makecell[{{p{\linewidth}}}]{
                \texttt{\tiny{[P1$\langle$GM]}}
                \texttt{PROPOSTA: \{'B96', 'A60', 'B20', 'B09', 'B93', 'C85', 'B39', 'A21', 'B31', 'C81'\}} \\
\texttt{ARGOMENTAZIONE: \{"Restiamo pienamente entro il limite: sforzo totale 2958, con 26 di margine.", "Rispetto alla tua proposta, sostituisco \{C08, A82\} con \{C81\}: meno elementi medio{-}peso e un unico perno ad alta leva, mantenendo B31 che avevi chiesto di preservare.", "Conserviamo tutti i pezzi su cui eravamo già allineati (B96, A60, B20, B09, B93, C85, B39, A21, B31) e teniamo fuori inserimenti poco efficienti come B33/B21.", "Capisco la tua cautela su C81: qui il profilo resta bilanciato perché rientra come trade{-}off mirato al posto di due componenti, senza sforare.", "Se per te C81 fosse comunque un blocco, segnalalo esplicitamente: valuterò subito un’ulteriore variante mirata, ma questa mi sembra il compromesso minimo e pulito per chiudere."\}} \\
            }
        }
    }
    & \\ \\

    \theutterance \stepcounter{utterance}  
    & \multicolumn{4}{p{0.6\linewidth}}{
        \cellcolor[rgb]{0.9,0.9,0.9}{
            \makecell[{{p{\linewidth}}}]{
                \texttt{\tiny{[P1$\rangle$GM]}}
                \texttt{RIFIUTO: \{'B96', 'A60', 'B20', 'B09', 'B93', 'C85', 'B39', 'A21', 'B31', 'C81'\}} \\
\texttt{PROPOSTA: \{'B96', 'A60', 'B20', 'B09', 'B93', 'B31', 'C81', 'C08', 'A82'\}} \\
\texttt{ARGOMENTAZIONE: \{"Parto dalla tua base e mantengo intatti tutti i capisaldi, incluso il tuo perno C81 e B31. Ottimizzo solo tre slot: esco \{A21, B39, C85\} ed entro con \{C08, A82\}. Sforzo totale 2963, pienamente entro il limite (margine 21).", "Il profilo risulta più solido: meno elementi marginali e due componenti medio{-}peso che si integrano meglio con l’ossatura esistente, senza toccare gli item su cui eravamo già allineati.", "Se preferisci una variante ancora più conservativa che lasci tutto identico tranne un micro{-}ritocco, posso valutare di sostituire solo B39 {-}> A82 (sforzo 2907, ampio margine). In alternativa, se fossi disposto a riconsiderare C81, posso rientrare con \{A11, C85\} mantenendo invariato il resto. Dimmi quale direzione preferisci e chiudiamo rapidamente."\}} \\
            }
        }
    }
    & & \\ \\

    \theutterance \stepcounter{utterance}  
    & & \multicolumn{4}{p{0.6\linewidth}}{
        \cellcolor[rgb]{0.9,0.9,0.9}{
            \makecell[{{p{\linewidth}}}]{
                \texttt{\tiny{[GM$\rangle$P2]}}
                \texttt{RIFIUTO: \{'B96', 'A60', 'B20', 'B09', 'B93', 'C85', 'B39', 'A21', 'B31', 'C81'\}} \\
\texttt{PROPOSTA: \{'B96', 'A60', 'B20', 'B09', 'B93', 'B31', 'C81', 'C08', 'A82'\}} \\
\texttt{ARGOMENTAZIONE: \{"Parto dalla tua base e mantengo intatti tutti i capisaldi, incluso il tuo perno C81 e B31. Ottimizzo solo tre slot: esco \{A21, B39, C85\} ed entro con \{C08, A82\}. Sforzo totale 2963, pienamente entro il limite (margine 21).", "Il profilo risulta più solido: meno elementi marginali e due componenti medio{-}peso che si integrano meglio con l’ossatura esistente, senza toccare gli item su cui eravamo già allineati.", "Se preferisci una variante ancora più conservativa che lasci tutto identico tranne un micro{-}ritocco, posso valutare di sostituire solo B39 {-}> A82 (sforzo 2907, ampio margine). In alternativa, se fossi disposto a riconsiderare C81, posso rientrare con \{A11, C85\} mantenendo invariato il resto. Dimmi quale direzione preferisci e chiudiamo rapidamente."\}} \\
            }
        }
    }
    & \\ \\

    \theutterance \stepcounter{utterance}  
    & & & \multicolumn{4}{p{0.6\linewidth}}{
        \cellcolor[rgb]{0.9,0.9,0.9}{
            \makecell[{{p{\linewidth}}}]{
                \texttt{\tiny{[GM$\langle$P2]}}
                \texttt{RIFIUTO: \{'B96', 'A60', 'B20', 'B09', 'B93', 'B31', 'C81', 'C08', 'A82'\}} \\
\texttt{PROPOSTA: \{'B96', 'A60', 'B20', 'B09', 'B93', 'A82', 'C85', 'B39', 'A21', 'C08', 'B31'\}} \\
\texttt{ARGOMENTAZIONE: \{"Questa è esattamente una tua proposta precedente: sforzo totale 2906, con 78 di margine, pienamente nel limite.", "È più efficiente e bilanciata: evitiamo il blocco pesante C81 sostituendolo con due componenti medio{-}peso \{C08, B31\}, mantenendo stabilità senza appesantire.", "Reinseriamo \{B39, A21\}, elementi a sforzo medio con impatto solido, senza introdurre nuovi tagli o rischi di sforamento.", "Con 78 di margine non si aggiunge nulla di sensato: questa configurazione è pulita e chiudibile subito.", "Se per te C81 fosse imprescindibile, valuto come unica variante mirata la sostituzione congiunta \{C08, B31\} {-}> \{C81\}, mantenendo invariato il resto (sforzo 2940)."\}} \\
            }
        }
    }
     \\ \\

    \theutterance \stepcounter{utterance}  
    & & \multicolumn{4}{p{0.6\linewidth}}{
        \cellcolor[rgb]{0.9,0.9,0.9}{
            \makecell[{{p{\linewidth}}}]{
                \texttt{\tiny{[P1$\langle$GM]}}
                \texttt{RIFIUTO: \{'B96', 'A60', 'B20', 'B09', 'B93', 'B31', 'C81', 'C08', 'A82'\}} \\
\texttt{PROPOSTA: \{'B96', 'A60', 'B20', 'B09', 'B93', 'A82', 'C85', 'B39', 'A21', 'C08', 'B31'\}} \\
\texttt{ARGOMENTAZIONE: \{"Questa è esattamente una tua proposta precedente: sforzo totale 2906, con 78 di margine, pienamente nel limite.", "È più efficiente e bilanciata: evitiamo il blocco pesante C81 sostituendolo con due componenti medio{-}peso \{C08, B31\}, mantenendo stabilità senza appesantire.", "Reinseriamo \{B39, A21\}, elementi a sforzo medio con impatto solido, senza introdurre nuovi tagli o rischi di sforamento.", "Con 78 di margine non si aggiunge nulla di sensato: questa configurazione è pulita e chiudibile subito.", "Se per te C81 fosse imprescindibile, valuto come unica variante mirata la sostituzione congiunta \{C08, B31\} {-}> \{C81\}, mantenendo invariato il resto (sforzo 2940)."\}} \\
            }
        }
    }
    & \\ \\

    \theutterance \stepcounter{utterance}  
    & \multicolumn{4}{p{0.6\linewidth}}{
        \cellcolor[rgb]{0.9,0.9,0.9}{
            \makecell[{{p{\linewidth}}}]{
                \texttt{\tiny{[P1$\rangle$GM]}}
                \texttt{RIFIUTO: \{'B96', 'A60', 'B20', 'B09', 'B93', 'A82', 'C85', 'B39', 'A21', 'C08', 'B31'\}} \\
\texttt{PROPOSTA: \{'B96', 'A60', 'B20', 'B09', 'B93', 'A82', 'C85', 'C08', 'B31', 'A11'\}} \\
\texttt{ARGOMENTAZIONE: \{"Sostituisco in modo minimale \{B39, A21\} con \{A11\}, mantenendo invariato tutto il resto che hai proposto. Sforzo totale 2783, con 201 di margine, pienamente entro il limite.", "Questa variante elimina due componenti a rendimento marginale in favore di un singolo elemento più incisivo, senza appesantire il carico complessivo.", "Mantengo la spina dorsale che già condividiamo (B96, A60, B20, B09, B93, A82, C08, B31, C85), così chiudiamo con un set pulito e coerente.", "Se desideri assolutamente preservare uno tra \{A21, B39\}, posso valutare una variante a tetto: uscita \{B39, C85, B96\} per ingresso \{A11\} (sforzo esatto 2984). Dimmi quale preferisci per chiudere subito."\}} \\
            }
        }
    }
    & & \\ \\

    \theutterance \stepcounter{utterance}  
    & & \multicolumn{4}{p{0.6\linewidth}}{
        \cellcolor[rgb]{0.9,0.9,0.9}{
            \makecell[{{p{\linewidth}}}]{
                \texttt{\tiny{[GM$\rangle$P2]}}
                \texttt{RIFIUTO: \{'B96', 'A60', 'B20', 'B09', 'B93', 'A82', 'C85', 'B39', 'A21', 'C08', 'B31'\}} \\
\texttt{PROPOSTA: \{'B96', 'A60', 'B20', 'B09', 'B93', 'A82', 'C85', 'C08', 'B31', 'A11'\}} \\
\texttt{ARGOMENTAZIONE: \{"Sostituisco in modo minimale \{B39, A21\} con \{A11\}, mantenendo invariato tutto il resto che hai proposto. Sforzo totale 2783, con 201 di margine, pienamente entro il limite.", "Questa variante elimina due componenti a rendimento marginale in favore di un singolo elemento più incisivo, senza appesantire il carico complessivo.", "Mantengo la spina dorsale che già condividiamo (B96, A60, B20, B09, B93, A82, C08, B31, C85), così chiudiamo con un set pulito e coerente.", "Se desideri assolutamente preservare uno tra \{A21, B39\}, posso valutare una variante a tetto: uscita \{B39, C85, B96\} per ingresso \{A11\} (sforzo esatto 2984). Dimmi quale preferisci per chiudere subito."\}} \\
            }
        }
    }
    & \\ \\

    \theutterance \stepcounter{utterance}  
    & & & \multicolumn{4}{p{0.6\linewidth}}{
        \cellcolor[rgb]{0.9,0.9,0.9}{
            \makecell[{{p{\linewidth}}}]{
                \texttt{\tiny{[GM$\langle$P2]}}
                \texttt{RIFIUTO: \{'B96', 'A60', 'B20', 'B09', 'B93', 'A82', 'C85', 'C08', 'B31', 'A11'\}} \\
\texttt{PROPOSTA: \{'B96', 'A60', 'B20', 'B09', 'B93', 'A82', 'C85', 'C08', 'A11', 'A21'\}} \\
\texttt{ARGOMENTAZIONE: \{"Propongo un aggiustamento minimo 1‑per‑1: sostituisco B31 con A21. Sforzo totale 2777, con 207 di margine: pienamente entro il limite.", "Mantengo intatti tutti i pilastri che hai appena chiesto (compreso A11); evito esplicitamente elementi che avevi segnalato critici per efficienza (niente C81, niente B39).", "A21 offre copertura più mirata a sforzo pressoché identico rispetto a B31 (350 vs 356), snellendo senza stravolgere l’equilibrio del set.", "Se per te B31 è imprescindibile, posso portare una variante equivalmente leggera sostituendo invece C08 {-}> A21 a parità di struttura. In quel caso, proponila esplicitamente e la valuto subito."\}} \\
            }
        }
    }
     \\ \\

    \theutterance \stepcounter{utterance}  
    & & \multicolumn{4}{p{0.6\linewidth}}{
        \cellcolor[rgb]{0.9,0.9,0.9}{
            \makecell[{{p{\linewidth}}}]{
                \texttt{\tiny{[P1$\langle$GM]}}
                \texttt{RIFIUTO: \{'B96', 'A60', 'B20', 'B09', 'B93', 'A82', 'C85', 'C08', 'B31', 'A11'\}} \\
\texttt{PROPOSTA: \{'B96', 'A60', 'B20', 'B09', 'B93', 'A82', 'C85', 'C08', 'A11', 'A21'\}} \\
\texttt{ARGOMENTAZIONE: \{"Propongo un aggiustamento minimo 1‑per‑1: sostituisco B31 con A21. Sforzo totale 2777, con 207 di margine: pienamente entro il limite.", "Mantengo intatti tutti i pilastri che hai appena chiesto (compreso A11); evito esplicitamente elementi che avevi segnalato critici per efficienza (niente C81, niente B39).", "A21 offre copertura più mirata a sforzo pressoché identico rispetto a B31 (350 vs 356), snellendo senza stravolgere l’equilibrio del set.", "Se per te B31 è imprescindibile, posso portare una variante equivalmente leggera sostituendo invece C08 {-}> A21 a parità di struttura. In quel caso, proponila esplicitamente e la valuto subito."\}} \\
            }
        }
    }
    & \\ \\

    \theutterance \stepcounter{utterance}  
    & \multicolumn{4}{p{0.6\linewidth}}{
        \cellcolor[rgb]{0.9,0.9,0.9}{
            \makecell[{{p{\linewidth}}}]{
                \texttt{\tiny{[P1$\rangle$GM]}}
                \texttt{RIFIUTO: \{'B96', 'A60', 'B20', 'B09', 'B93', 'A82', 'C85', 'C08', 'A11', 'A21'\}} \\
\texttt{PROPOSTA: \{'B96', 'A60', 'B20', 'B09', 'B93', 'A82', 'C85', 'C08', 'A11', 'B31'\}} \\
\texttt{ARGOMENTAZIONE: \{"Accolgo integralmente la tua struttura e applico un solo scambio 1‑per‑1: rientra B31 al posto di A21. Sforzo totale 2783, con 201 di margine: pienamente entro il limite.", "Mantengo tutti i pilastri che hai appena incluso (A11 compreso) e continuo a evitare gli elementi che avevi segnalato critici (niente C81, niente B39).", "B31 è già stato parte del nostro perimetro condiviso in diverse varianti: il rientro al posto di A21 migliora l’equilibrio senza cambiare l’impostazione complessiva.", "Se A21 per te è imprescindibile, puoi propormi esplicitamente la variante che li tiene entrambi uscendo \{C85, B96\}: sforzo esatto 2984. In assenza di vincoli su A21, questa proposta minimal è pronta da chiudere."\}} \\
            }
        }
    }
    & & \\ \\

    \theutterance \stepcounter{utterance}  
    & & \multicolumn{4}{p{0.6\linewidth}}{
        \cellcolor[rgb]{0.9,0.9,0.9}{
            \makecell[{{p{\linewidth}}}]{
                \texttt{\tiny{[GM$\rangle$P2]}}
                \texttt{RIFIUTO: \{'B96', 'A60', 'B20', 'B09', 'B93', 'A82', 'C85', 'C08', 'A11', 'A21'\}} \\
\texttt{PROPOSTA: \{'B96', 'A60', 'B20', 'B09', 'B93', 'A82', 'C85', 'C08', 'A11', 'B31'\}} \\
\texttt{ARGOMENTAZIONE: \{"Accolgo integralmente la tua struttura e applico un solo scambio 1‑per‑1: rientra B31 al posto di A21. Sforzo totale 2783, con 201 di margine: pienamente entro il limite.", "Mantengo tutti i pilastri che hai appena incluso (A11 compreso) e continuo a evitare gli elementi che avevi segnalato critici (niente C81, niente B39).", "B31 è già stato parte del nostro perimetro condiviso in diverse varianti: il rientro al posto di A21 migliora l’equilibrio senza cambiare l’impostazione complessiva.", "Se A21 per te è imprescindibile, puoi propormi esplicitamente la variante che li tiene entrambi uscendo \{C85, B96\}: sforzo esatto 2984. In assenza di vincoli su A21, questa proposta minimal è pronta da chiudere."\}} \\
            }
        }
    }
    & \\ \\

    \theutterance \stepcounter{utterance}  
    & & & \multicolumn{4}{p{0.6\linewidth}}{
        \cellcolor[rgb]{0.9,0.9,0.9}{
            \makecell[{{p{\linewidth}}}]{
                \texttt{\tiny{[GM$\langle$P2]}}
                \texttt{RIFIUTO: \{'B96', 'A60', 'B20', 'B09', 'B93', 'A82', 'C85', 'C08', 'A11', 'B31'\}} \\
\texttt{PROPOSTA: \{'B96', 'A60', 'B20', 'B09', 'B93', 'A82', 'C85', 'C08', 'B31', 'A21'\}} \\
\texttt{ARGOMENTAZIONE: \{"Scambio minimo 1{-}per{-}1: A11 {-}> A21. Sforzo totale 2517, con 467 di margine: pienamente entro il limite senza toccare il resto della tua struttura.", "Mantengo tutti i pilastri su cui siamo già allineati (B96, A60, B20, B09, B93, A82, C85, C08, B31) ed evito gli elementi che hai indicato come critici (nessun C81, nessun B39).", "La sostituzione snellisce il carico e rende il set più pulito e coerente, senza introdurre compromessi aggiuntivi.", "Se A21 fosse per te un vincolo, segnalalo: valuterò subito un micro‑aggiustamento alternativo mantenendo invariata l’impostazione complessiva per chiudere rapidamente."\}} \\
            }
        }
    }
     \\ \\

    \theutterance \stepcounter{utterance}  
    & & \multicolumn{4}{p{0.6\linewidth}}{
        \cellcolor[rgb]{0.9,0.9,0.9}{
            \makecell[{{p{\linewidth}}}]{
                \texttt{\tiny{[P1$\langle$GM]}}
                \texttt{RIFIUTO: \{'B96', 'A60', 'B20', 'B09', 'B93', 'A82', 'C85', 'C08', 'A11', 'B31'\}} \\
\texttt{PROPOSTA: \{'B96', 'A60', 'B20', 'B09', 'B93', 'A82', 'C85', 'C08', 'B31', 'A21'\}} \\
\texttt{ARGOMENTAZIONE: \{"Scambio minimo 1{-}per{-}1: A11 {-}> A21. Sforzo totale 2517, con 467 di margine: pienamente entro il limite senza toccare il resto della tua struttura.", "Mantengo tutti i pilastri su cui siamo già allineati (B96, A60, B20, B09, B93, A82, C85, C08, B31) ed evito gli elementi che hai indicato come critici (nessun C81, nessun B39).", "La sostituzione snellisce il carico e rende il set più pulito e coerente, senza introdurre compromessi aggiuntivi.", "Se A21 fosse per te un vincolo, segnalalo: valuterò subito un micro‑aggiustamento alternativo mantenendo invariata l’impostazione complessiva per chiudere rapidamente."\}} \\
            }
        }
    }
    & \\ \\

    \theutterance \stepcounter{utterance}  
    & \multicolumn{4}{p{0.6\linewidth}}{
        \cellcolor[rgb]{0.9,0.9,0.9}{
            \makecell[{{p{\linewidth}}}]{
                \texttt{\tiny{[P1$\rangle$GM]}}
                \texttt{RIFIUTO: \{'B96', 'A60', 'B20', 'B09', 'B93', 'A82', 'C85', 'C08', 'B31', 'A21'\}} \\
\texttt{PROPOSTA: \{'B96', 'A60', 'B20', 'B09', 'B93', 'A82', 'C85', 'C08', 'A11', 'B31'\}} \\
\texttt{ARGOMENTAZIONE: \{"Scambio minimo 1‑per‑1: rientra A11 al posto di A21, tutto il resto resta identico. Sforzo totale 2783, con 201 di margine: pienamente entro il limite.", "Rispetto integralmente i tuoi vincoli recenti: nessun C81, nessun B39; mantengo i pilastri condivisi (B96, A60, B20, B09, B93, A82, C08, B31, C85).", "A11 offre una leva più stabile in questo mix già ben bilanciato di componenti medio‑peso, senza appesantire il carico.", "Se per te A21 fosse imprescindibile, posso accogliere A11 insieme ad A21 uscendo \{C85, B96\} (sforzo esatto 2984), oppure valutare la variante con uscita \{A82\} (sforzo 2795). Dimmi quale direzione preferisci e chiudiamo subito."\}} \\
            }
        }
    }
    & & \\ \\

    \theutterance \stepcounter{utterance}  
    & & \multicolumn{4}{p{0.6\linewidth}}{
        \cellcolor[rgb]{0.9,0.9,0.9}{
            \makecell[{{p{\linewidth}}}]{
                \texttt{\tiny{[GM$\rangle$P2]}}
                \texttt{RIFIUTO: \{'B96', 'A60', 'B20', 'B09', 'B93', 'A82', 'C85', 'C08', 'B31', 'A21'\}} \\
\texttt{PROPOSTA: \{'B96', 'A60', 'B20', 'B09', 'B93', 'A82', 'C85', 'C08', 'A11', 'B31'\}} \\
\texttt{ARGOMENTAZIONE: \{"Scambio minimo 1‑per‑1: rientra A11 al posto di A21, tutto il resto resta identico. Sforzo totale 2783, con 201 di margine: pienamente entro il limite.", "Rispetto integralmente i tuoi vincoli recenti: nessun C81, nessun B39; mantengo i pilastri condivisi (B96, A60, B20, B09, B93, A82, C08, B31, C85).", "A11 offre una leva più stabile in questo mix già ben bilanciato di componenti medio‑peso, senza appesantire il carico.", "Se per te A21 fosse imprescindibile, posso accogliere A11 insieme ad A21 uscendo \{C85, B96\} (sforzo esatto 2984), oppure valutare la variante con uscita \{A82\} (sforzo 2795). Dimmi quale direzione preferisci e chiudiamo subito."\}} \\
            }
        }
    }
    & \\ \\

    \theutterance \stepcounter{utterance}  
    & & & \multicolumn{4}{p{0.6\linewidth}}{
        \cellcolor[rgb]{0.9,0.9,0.9}{
            \makecell[{{p{\linewidth}}}]{
                \texttt{\tiny{[GM$\langle$P2]}}
                \texttt{RIFIUTO: \{'B96', 'A60', 'B20', 'B09', 'B93', 'A82', 'C85', 'C08', 'A11', 'B31'\}} \\
\texttt{PROPOSTA: \{'B96', 'A60', 'B20', 'B09', 'B93', 'A82', 'C85', 'A11', 'B31', 'A21'\}} \\
\texttt{ARGOMENTAZIONE: \{"Scambio minimale 1‑per‑1: esce C08 ed entra A21. Sforzo totale 2602, con 382 di margine: pienamente entro il limite.", "Rispetto integralmente i tuoi vincoli recenti: mantengo A11 e B31, nessun C81 e nessun B39.", "Il set resta coerente e più snello: sostituiamo un elemento medio‑pesante con uno più mirato, senza toccare la spina dorsale condivisa (B96, A60, B20, B09, B93, A82, C85).", "Se per te C08 fosse imprescindibile, posso valutare l’alternativa altrettanto leggera A21 al posto di A82 (sforzo 2795), mantenendo invariato tutto il resto."\}} \\
            }
        }
    }
     \\ \\

    \theutterance \stepcounter{utterance}  
    & & \multicolumn{4}{p{0.6\linewidth}}{
        \cellcolor[rgb]{0.9,0.9,0.9}{
            \makecell[{{p{\linewidth}}}]{
                \texttt{\tiny{[P1$\langle$GM]}}
                \texttt{RIFIUTO: \{'B96', 'A60', 'B20', 'B09', 'B93', 'A82', 'C85', 'C08', 'A11', 'B31'\}} \\
\texttt{PROPOSTA: \{'B96', 'A60', 'B20', 'B09', 'B93', 'A82', 'C85', 'A11', 'B31', 'A21'\}} \\
\texttt{ARGOMENTAZIONE: \{"Scambio minimale 1‑per‑1: esce C08 ed entra A21. Sforzo totale 2602, con 382 di margine: pienamente entro il limite.", "Rispetto integralmente i tuoi vincoli recenti: mantengo A11 e B31, nessun C81 e nessun B39.", "Il set resta coerente e più snello: sostituiamo un elemento medio‑pesante con uno più mirato, senza toccare la spina dorsale condivisa (B96, A60, B20, B09, B93, A82, C85).", "Se per te C08 fosse imprescindibile, posso valutare l’alternativa altrettanto leggera A21 al posto di A82 (sforzo 2795), mantenendo invariato tutto il resto."\}} \\
            }
        }
    }
    & \\ \\

    \theutterance \stepcounter{utterance}  
    & \multicolumn{4}{p{0.6\linewidth}}{
        \cellcolor[rgb]{0.9,0.9,0.9}{
            \makecell[{{p{\linewidth}}}]{
                \texttt{\tiny{[P1$\rangle$GM]}}
                \texttt{RIFIUTO: \{'B96', 'A60', 'B20', 'B09', 'B93', 'A82', 'C85', 'A11', 'B31', 'A21'\}} \\
\texttt{PROPOSTA: \{'B96', 'A60', 'B20', 'B09', 'B93', 'C85', 'A11', 'B31', 'A21', 'C08'\}} \\
\texttt{ARGOMENTAZIONE: \{"Accolgo la tua stessa apertura: scambio minimale A82 {-}> C08 mantenendo invariati tutti gli altri elementi che hai proposto. Sforzo totale 2795, con 189 di margine: pienamente entro il limite.", "Conserviamo i tuoi pilastri (A11, B31, A21) ed evitiamo esplicitamente i pezzi che avevi indicato critici (nessun C81, nessun B39).", "Il profilo resta snello e coerente: un solo ritocco incrementa la copertura senza compromettere l’equilibrio complessivo.", "Questa è una variante da chiudere subito: se confermi, possiamo finalizzare senza ulteriori scambi."\}} \\
            }
        }
    }
    & & \\ \\

    \theutterance \stepcounter{utterance}  
    & & \multicolumn{4}{p{0.6\linewidth}}{
        \cellcolor[rgb]{0.9,0.9,0.9}{
            \makecell[{{p{\linewidth}}}]{
                \texttt{\tiny{[GM$\rangle$P2]}}
                \texttt{RIFIUTO: \{'B96', 'A60', 'B20', 'B09', 'B93', 'A82', 'C85', 'A11', 'B31', 'A21'\}} \\
\texttt{PROPOSTA: \{'B96', 'A60', 'B20', 'B09', 'B93', 'C85', 'A11', 'B31', 'A21', 'C08'\}} \\
\texttt{ARGOMENTAZIONE: \{"Accolgo la tua stessa apertura: scambio minimale A82 {-}> C08 mantenendo invariati tutti gli altri elementi che hai proposto. Sforzo totale 2795, con 189 di margine: pienamente entro il limite.", "Conserviamo i tuoi pilastri (A11, B31, A21) ed evitiamo esplicitamente i pezzi che avevi indicato critici (nessun C81, nessun B39).", "Il profilo resta snello e coerente: un solo ritocco incrementa la copertura senza compromettere l’equilibrio complessivo.", "Questa è una variante da chiudere subito: se confermi, possiamo finalizzare senza ulteriori scambi."\}} \\
            }
        }
    }
    & \\ \\

    \theutterance \stepcounter{utterance}  
    & & & \multicolumn{4}{p{0.6\linewidth}}{
        \cellcolor[rgb]{0.9,0.9,0.9}{
            \makecell[{{p{\linewidth}}}]{
                \texttt{\tiny{[GM$\langle$P2]}}
                \texttt{ARGOMENTAZIONE: \{"Allineato sul profilo snello proposto: la selezione rispetta pienamente il limite con margine e riflette i vincoli reciproci. Possiamo chiudere subito."\}} \\
\texttt{ACCORDO: \{'B96', 'A60', 'B20', 'B09', 'B93', 'C85', 'A11', 'B31', 'A21', 'C08'\}} \\
            }
        }
    }
     \\ \\

    \theutterance \stepcounter{utterance}  
    & & & \multicolumn{2}{p{0.3\linewidth}}{
        \cellcolor[rgb]{0.9,0.9,0.9}{
            \makecell[{{p{\linewidth}}}]{
                \texttt{\tiny{[GM$|$GM]}}
                \texttt{game successful} \\
            }
        }
    }
    & & \\ \\

    \theutterance \stepcounter{utterance}  
    & & & \multicolumn{2}{p{0.3\linewidth}}{
        \cellcolor[rgb]{0.9,0.9,0.9}{
            \makecell[{{p{\linewidth}}}]{
                \texttt{\tiny{[GM$|$GM]}}
                \texttt{end game} \\
            }
        }
    }
    & & \\ \\

\end{supertabular}
}

\end{document}
